\begin{table}[tp]
\centering
\caption{Comparison for \href{https://github.com/k3nfalt/longtails/blob/main/LongTails/Setups/SpikedSubGaussianConvex.lean}{SpikedSubGaussianConvex}.}\label{tab:SpikedSubGaussianConvex}
\parbox{\linewidth}{\footnotesize Convex, globally $G$-Lipschitz $f:\mathbb R^d\to\mathbb R$ has a global minimizer $\|x^\star\|\le R$; $G,R,\varepsilon>0$. Fresh IID samples from a probability measure $p$ feed a jointly measurable, integrable oracle $g$, with $\int g(x,\omega)\,dp=v(x)\in\partial f(x)$. For every $x,u,t$, the exponential moment is integrable and at most $\exp(t^2[v_{\rm rest}\|u\|^2+(v_{\rm top}-v_{\rm rest})u_0^2]/2)$ for $\langle g-v,u\rangle$, where $0\le v_{\rm rest}\le v_{\rm top}$ and the spike direction is known. Target: measurable $\hat x$ with $P\{f(\hat x)-f(x^\star)>\varepsilon\}\le1/3$.}\par\smallskip
\begin{threeparttable}
\footnotesize
\setlength{\tabcolsep}{4pt}\renewcommand{\arraystretch}{1.25}
\begin{tabularx}{\linewidth}{@{}>{\raggedright\arraybackslash}p{0.20\linewidth}>{\raggedright\arraybackslash}p{0.32\linewidth}>{\raggedright\arraybackslash}p{0.25\linewidth}>{\raggedright\arraybackslash}X@{}}
\toprule
\textbf{Result} & \textbf{Oracle complexity} & \textbf{Novelty} & \textbf{Limitations} \\
\midrule
\textbf{(1)} \href{https://github.com/k3nfalt/longtails/blob/main/LongTails/Theorems/SpikedSubGaussianSoftCutConvex.lean}{Spiked\allowbreak Sub\allowbreak Gaussian\allowbreak Soft\allowbreak Cut\allowbreak Convex} \textbf{(Ours)} $(\star\star)$ \href{https://github.com/k3nfalt/longtails/blob/main/scripts/reviews/SpikedSubGaussianSoftCutConvex.md}{(review)} & $\mathcal{O}\!\left((d+\tau R^2/\varepsilon^2)\ell^5\right)$ & & \\
\midrule[1.2pt]
\textbf{(2)} ESGO cutting planes \citep{spikedMOSZ2025v1}, Cor.~3 & $\mathcal{O}\!\left((d+d v_{\rm top}R^2/\varepsilon^2)\ell^{\mathcal{O}(1)}\right)$ & \textbf{(1)}: noise coefficient $\tau$ instead of $dv_{\rm top}$; gain up to $d$ when anisotropy dominates logarithms. For $G=R=1$, $\varepsilon=v_{\rm top}=d^{-1}$, $v_{\rm rest}=d^{-3}$: $\mathcal{O}(d\log^5d)$ versus $\mathcal{O}(d^2\log^{\mathcal{O}(1)}d)$. & \textbf{(1)} needs a known spike and sub-Gaussian MGF; \textbf{(2)} allows directional sub-exponential tails. Isotropic noise gives no polynomial gain. \\
\midrule
\textbf{(3)} SGD, derived probability bound \citep{spikedMOSZ2025v1}, \S1 & $\mathcal{O}\!\left(1+(G^2+\tau)R^2/\varepsilon^2\right)$ & \textbf{(1)} removes $G^2R^2/\varepsilon^2$, adding $d\ell^5$; wins if $G^2R^2/\varepsilon^2\gg\max\{d,\tau R^2/\varepsilon^2\}\ell^5$. In \textbf{(2)}'s example: $\mathcal{O}(d\log^5d)$ versus $\mathcal{O}(d^2)$. & \textbf{(1)} requires known-spike sub-Gaussian noise; \textbf{(3)} needs only finite variance and avoids $\ell^5$. \textbf{(3)} is cheaper when $G^2R^2/\varepsilon^2\lesssim d$. \\
\bottomrule
\end{tabularx}
\begin{tablenotes}[flushleft]\footnotesize
\item[] \textbf{Abbr:} $d=m+1$, $\tau=v_{\rm top}+(d-1)v_{\rm rest}$, $\ell=1+\log(d(1+GR/\varepsilon))$. Costs count individual stochastic subgradient calls, including batches. $\ell^{\mathcal{O}(1)}$ denotes the polylogarithmic factors suppressed in the cited bound. ESGO: sub-exponential stochastic gradient oracle, here with scale $\sigma_E\asymp\sqrt{dv_{\rm top}}$. SGD: stochastic gradient descent; its expected-error bound uses $\mathbb E\|g\|^2\le G^2+\tau$, converted by Markov's inequality.
\end{tablenotes}
\end{threeparttable}
\end{table}
\clearpage
